\originalchapter{8.06 作业2：高阶微扰与环上极化}
来源为Spring 2018 Assignment 2, 2018年2月14日发布, Barton Zwiebach. 本轮第1、2题完整中文转述与\textbf{AI独立解答}. 第3题署Schiff改编、第4题外书Feynman--Hellmann题、第5题署Griffiths变题, 各具体表达单核, 不整份禁入, 未计入本轮完成数. 一般微扰和氢Stark主题已有本册独立推导, 不冒称取得商业教材.
\originalsection{第1题：高阶非简并微扰与归一化}
\begin{exercise}
$H=H_0+\lambda V$, 非简并$\ket n$, 采用中间归一$\braket n{n_\lambda}=1$. 求二阶态修正和三阶能修正. 展$Z_n$至$\lambda^2$, $Z_n^{-1}=\braket{n_\lambda}{n_\lambda}$, 新态测到旧$n$的概率是什么?
\end{exercise}
\begin{solution}
记$D_k=E_n^0-E_k^0$, $k\ne n$, 一阶态$c_k^{(1)}=V_{kn}/D_k$, $E^{(1)}=V_{nn}$. 将本征方程$\lambda^2$项投影$k\ne n$, 得$D_kc_k^{(2)}=\sum_{l\ne n}V_{kl}V_{ln}/D_l-V_{nn}V_{kn}/D_k$. 所以
\[
\ket{n^{(2)}}=\sum_{k\ne n}\left[\sum_{l\ne n}\frac{V_{kl}V_{ln}}{D_kD_l}-\frac{V_{nn}V_{kn}}{D_k^2}\right]\ket k.
\]
投影$n$的三阶方程给
\[
E^{(3)}=\sum_{k,l\ne n}\frac{V_{nk}V_{kl}V_{ln}}{D_kD_l}
-V_{nn}\sum_{k\ne n}\frac{|V_{kn}|^2}{D_k^2}.
\]
Hermitian时为实, 第一和的共轭可交换$k,l$核验. 中间归一各阶态与$\ket n$正交, 故范数$1+\lambda^2\sum|V_{kn}|^2/D_k^2+O(\lambda^3)$, $Z_n=1-\lambda^2\sum|V_{kn}|^2/D_k^2+O(\lambda^3)$. 物理单位态$\sqrt{Z_n}\ket{n_\lambda}$测旧$n$概率恰$Z_n$（中间归一时振幅旧分量1）, 不是未经归一就断言概率1. 若只保留一阶态再平方, 得相同至二阶结果.
\end{solution}
\originalsection{第2题：简并环态的二阶极化}
\begin{exercise}
环半径$a$, $\psi_n=\e^{\ii n\phi}/\sqrt{2\pi}$, $E_n=E_1n^2$, $E_1=\hbar^2/(2ma^2)$. 电场$\epsilon$沿$x$, 扰动$V=-q\epsilon a\cos\phi$. 求$\ket{\pm1}$简并能级的二阶修正及零阶好态.
\end{exercise}
\begin{solution}
令$v=q\epsilon a/2$, 矩阵元$V_{mn}=-v(\delta_{m,n+1}+\delta_{m,n-1})$, 简并内部一阶矩阵0. 外部中间态为$0,\pm2$, 对角二阶$v^2(1/E_1+1/(E_1-4E_1))=2v^2/(3E_1)$. 两态间只共用中间$0$, 非对角$v^2/E_1$. 所以二阶有效矩阵为$(v^2/E_1)\begin{pmatrix}2/3&1\\1&2/3\end{pmatrix}$. 好态$(\ket1+\ket{-1})/\sqrt2$为$\cos\phi/\sqrt\pi$, 位移$5v^2/(3E_1)$; $(\ket1-\ket{-1})/\sqrt2$到整体相位为$\sin\phi/\sqrt\pi$, 位移$-v^2/(3E_1)$. 对零阶简并态直接分别套非简并二阶公式会漏非对角耦合而得相同错位移.
\end{solution}
